\cleardoublepage

\chapter*{Resumen\markboth{Resumen}{Resumen}}

La seguridad en los sistemas modernos es un área en constante evolución, donde la integridad de los datos y el control de su funcionamiento juegan un papel fundamental.
Esto es debido a la presencia de amenazas y agentes externos que trabajan con el fin malicioso de comprometer dicha información, haciendo necesario el desarrollo de mecanismos de protección cada vez más robustos.\\

En este ámbito, una de las amenazas más críticas para los sistemas actualmente son los \textit{rootkits}, software malicioso diseñado para ocultarse en el núcleo del sistema (\textit{kernel}) y controlar su funcionamiento.
Al infiltrarse en los niveles más profundos, este \textit{malware} adquiere permisos de administrador, permitiéndole controlar la ejecución de procesos de forma arbitraria y manipular respuestas e información del sistema operativo, 
impidiendo obtener una visión verídica del estado real de los archivos y procesos del sistema.
Esto provoca una falta de integridad en la monitorización y control del mismo, donde por defecto se deposita una confianza implícita en una fuente de datos sensibles que puede estar comprometida ante dicho tipo de amenaza.\\

Para abordar esta problemática, en este proyecto se ha diseñado e implementado un \textit{Módulo del Kernel} (\textit{LKM}, \textit{Loadable Kernel Module}) en un sistema \textit{Linux} con el objetivo de protegerlo ante este tipo de amenazas.
Esta solución introduce un mecanismo de autenticación criptográfica basado en HMAC (\textit{Hash-based Message Authentication Code}) del contenido proporcionado por el \textit{kernel},
mediante una clave simétrica secreta que solo es conocida por el administrador, garantizando la integridad de la información ofrecida.\\

Los experimentos realizados aproximan un entorno real de operación donde la información sensible es solicitada desde el \textit{espacio de usuario} hacia el \textit{espacio del kernel}.
Con este fin, se define el contenido sensible del sistema a autenticar mediante un listado de todos los procesos en ejecución y sus respectivos metadatos, generado en el propio \textit{Módulo del Kernel}, sirviendo como ejemplo de información crítica que puede ser manipulada por un \textit{rootkit}.
Sobre este contenido, se calcula el HMAC con la respectiva clave secreta y se adjunta a la información proporcionada al usuario, de forma que este pueda verificar su integridad y autenticidad.\\

Para validar esta propuesta, se han desarrollado componentes de usuario que prueban la funcionalidad. Un primer componente genera un fichero cuyo descriptor es transferido al \textit{kernel} para que este lo referencie, estableciendo así un canal de comunicación seguro.
Por otro lado, la integridad de los datos reportados se confirma mediante un segundo componente de verificación que analiza la información escrita por el \textit{kernel}, calculando localmente el HMAC del contenido con la misma clave secreta y
contrastándolo posteriormente con el recibido para asegurar que coinciden y que el contenido es verídico.\\